\begin{claim}{D.25}{Proved}{All-prime dihedral compatibility and the multiplier-one obstruction}
The inverse limit $D$ over all the dihedral levels of D.24 has the
following explicit isomorphism of profinite groups:
\[
 D\xrightarrow{\ \sim\ }
 \left(\prod_{p\text{ odd}}\operatorname{Aff}(\Z_p)\right)
 \times H_2.                                                   \tag{D.25a}
\]
There is a canonical continuous homomorphism $GT_0\to D$.
The group $D$ contains an element of every prime order. In
particular, imposing order three throughout this entire
dihedral system still permits a multiplier whose $7$-adic
coordinate is a primitive cube root and whose other prime
coordinates are $1$.

Put $F=\widehat F_2=\langle x,y\rangle$ and define the continuous
epimorphism and its closed kernel by
\[
 \kappa:F\longrightarrow\prod_pF_2^{(p)},\qquad
 R=\ker\kappa.
\]
Write $\lambda:GT_0\to\widehat\Z^\times$ for the multiplier.
For every integer $d\geq2$, the elements $a\in\ker\lambda$
with $a^d=1$ correspond exactly to the following elements
$f\in R$. The continuous endomorphism
\[
 \alpha_f:F\to F,\qquad x\mapsto x,\quad y\mapsto f^{-1}yf
\]
must be an automorphism and satisfy
\[
 \theta(f)f=1,\qquad
 \tau^2(f)\tau(f)f=1,\qquad
 \prod_{j=0}^{d-1}\alpha_f^{j}(f)=1,                         \tag{D.25b}
\]
where the product is in increasing order of $j$,
$\theta(x)=y$, $\theta(y)=x$, $\tau(x)=y$, and
$\tau(y)=(xy)^{-1}$. Every such finite-order element has identity
image in $D$. The kernel $R$ itself maps onto the explicit
perfect group $A_5$ used in D.24.
The finite-order solutions of (D.25b) in the stated
automorphism domain are exactly $f=1$, by D.32's subsequent
full fixed-subgroup argument.  Thus the nontrivial prime-order
points of the dihedral inverse limit exhibited here do not
lift with that order to $GT_0$ when the prime is odd.
\end{claim}
\begin{proof}
It suffices to use the even levels $n=2^\alpha n_0$ with
$\alpha\geq1$ and $n_0$ odd; every odd level agrees with its
double. Apply the Chinese remainder isomorphism separately to
the translation coordinate modulo $n/2$ and to the multiplier
coordinate modulo $2n$ in (D.24b). The odd factors give
$\operatorname{Aff}(\Z/n_0)$, and the powers of two give
the finite groups defining $H_2$. These coordinate maps commute
with every reduction $n\mid q$. Conversely, specified odd
coordinates and a specified point of $H_2$ reconstruct every
finite pair $(k,u)$ by the Chinese remainder theorem. They
satisfy the parity condition because that condition involves
only the power-of-two coordinates. The resulting inverse
maps are continuous, being continuous at every finite
coordinate. This proves (D.25a), including its topology.
Restriction of the full inverse-limit point in D.23 defines
the indicated map from $GT_0$; surjectivity onto $D$ is not
used or asserted.

Let $q$ be a prime, and choose a prime divisor $p$ of $2^q-1$.
The residue of $2$ in $\F_p^\times$ has order $q$: its order
divides $q$, and order one would give $p\mid2-1$. Hence
$q\mid p-1$, in particular $p$ is odd and $p\ne q$.
The element $2\bmod p$ lifts to a unique $q$th root
$\zeta\in\Z_p^\times$. Explicitly, after a root $a$ modulo
$p^j$ is chosen, substitute $a+p^jt$, $0\leq t<p$,
in $X^q-1$ modulo $p^{j+1}$. Its coefficient of $t$ is
$p^jqa^{q-1}$, and $qa^{q-1}$ is a unit modulo $p$;
there is exactly one permitted value of $t$.
The compatible limit has order exactly $q$ because its
residue already does. Put $(0,\zeta)$ in the $p$-factor
of (D.25a), and put the identity in all the other factors.
This gives the asserted element of order $q$ in $D$.

For $q=3$, take $p=7$. The executed program
\texttt{code/part\_a/gt\_shadow\_all\_prime.py} constructs ten
successive roots of $X^2+X+1$ starting with $2\bmod7$, and
uses the exact Chinese remainder formulas to construct
compatible shadow coordinates at sixteen levels through
$n=2744$. At $n=14$ its affine pair is $(0,9\bmod28)$;
at $n=98$ it is $(0,177\bmod196)$. Their orders are three
and the latter reduces to the former. The preceding
root-lifting proof constructs the infinite point, while
the program verifies the displayed finite coordinates.
Thus the equations tested over all dihedral levels admit
prime-dependent torsion multipliers. They do not force
the global alternatives $1$ and $-1$.

We next prove the statement about $R$. Each projection
$F\to F_2^{(p)}$ is the canonical maximal pro-$p$ quotient
map. Their product is onto: in any finite collection of
finite $p$-group quotients at distinct primes, its image
surjects to each factor. Its order is consequently divisible
by every factor order, hence by their product, since these
orders are pairwise coprime. The image is the full product.
These finite products form the finite quotients of
$\prod_pF_2^{(p)}$, so the compact image of $F$ is the whole
product. Also $R\subset F'$, since the abelianization
$F^{\mathrm{ab}}=\widehat\Z^2$ is detected by the
abelianizations of all these pro-$p$ quotients.

Let $a=(1,f)\in GT_0$ have finite order dividing $d$.
It induces a finite-order automorphism of each free pro-$p$
quotient with trivial action on its abelianization. The
Magnus leading-term argument in D.10, which also applies
at $p=2$, makes this automorphism the identity. Its normalized
formula then gives
$f_p\in C_{F_2^{(p)}}(y)\cap(F_2^{(p)})'=1$ by A.7 and
the primitive abelianization coordinate of $y$.
Thus $f\in R$.

The two first equations of (D.25b) are precisely the defining
cycle equations at multiplier one; the required invertibility
is stated explicitly. To obtain the last equation, put
$w_d=f\alpha_f(f)\cdots\alpha_f^{d-1}(f)$.
Induction on $d$ gives
\[
 \alpha_f^d(x)=x,\qquad
 \alpha_f^d(y)=w_d^{-1}yw_d.
\]
Therefore $\alpha_f^d=1$ is equivalent to
$w_d\in C_F(y)$. But $w_d\in F'$ and
$C_F(y)\cap F'=\overline{\langle y\rangle}\cap F'=1$,
again by A.7 and abelianization. This proves the necessity
of the last equation and its sufficiency. The normalized
identification of $GT_0$ with these automorphisms gives
the asserted bijection. The kernel of the multiplier on
$D$ is a subgroup of the additive translation group
$\prod_p\Z_p$, which is torsion-free. Consequently every
finite-order multiplier-one element maps to the identity
of $D$.

Finally let $\eta:F\to A_5$ be the epimorphism from D.24.
The quotient $A_5/\eta(R)$ is a finite quotient of
$F/R=\prod_pF_2^{(p)}$ and is therefore nilpotent.
It is also perfect, since the executed permutation
calculation in D.24 proves $A_5'=A_5$. A finite nilpotent
perfect group is trivial: its lower central series would
otherwise be constantly equal to the nontrivial group
and could not terminate. Hence $\eta(R)=A_5$.
This identifies a nontrivial part of the residual group in the
exact equation (D.25b). The full three-cusp argument in D.32
proves that this equation, together with its stated cycle and
invertibility requirements, has only the identity solution.
The surjection $R\to A_5$ and the compatible points of $D$
remain the exact quotient calculations proved above.
\end{proof}
